% TOP_PROBLEM: {"id": 3299, "title": "Infinite uniquely hamiltonian graphs", "category_id": 3, "category": {"id": 3, "name": "graph_theory", "display_name": "Graph Theory", "description": "Problems involving graphs, networks, and their properties.", "slug": "graph-theory", "order_index": 3, "created_at": "2026-07-31T15:26:25.671Z"}, "status": "open", "problem_number": "OPG-484", "set_id": 12}
% FIRST_POSTED: 2026-10-01
% ENTRY_KIND: statement_only
% UnsolvedMath ID 3299; source catalogue code OPG-484
% !TeX program = lualatex
% Definition and sources only; this file is not an LLM solution attempt.
\documentclass[11pt,a4paper]{article}
\usepackage[margin=25mm]{geometry}
\usepackage{fontspec}
\setmainfont{lmroman10-regular.otf}[BoldFont=lmroman10-bold.otf,
 ItalicFont=lmroman10-italic.otf,BoldItalicFont=lmroman10-bolditalic.otf]
\usepackage{amsmath,amssymb,mathtools}
\usepackage{unicode-math}
\setmathfont{latinmodern-math.otf}
\AtBeginDocument{\let\setminus\smallsetminus}
\usepackage{xurl,hyperref}
\hypersetup{colorlinks=true,urlcolor=blue}
\setcounter{secnumdepth}{0}
\setlength{\parindent}{0pt}
\setlength{\parskip}{5pt}
\setlength{\emergencystretch}{3em}
\begin{document}
\section*{Infinite uniquely hamiltonian graphs}\label{3299}
{\small UnsolvedMath ID 3299 · OPG-484 · Graph Theory · OpenGarden}
\subsection{Definitions and mathematical statement}
Let $G$ be an infinite connected locally finite graph, meaning that each vertex has finite degree. A ray is a one-way infinite simple path. Two rays belong to the same end when, after deleting any finite vertex set, their tails lie in the same remaining component. The graph is one-ended if all its rays are equivalent in this sense.

The Freudenthal compactification $|G|$ adds one point for each end to the usual geometric realization of the graph by unit intervals. A neighborhood of an end contains the component carrying its rays after some finite vertex deletion, all ends carried by that component, and the adjacent open edge pieces entering it. A Hamilton circle is a subspace of $|G|$ homeomorphic to the ordinary circle $S^1$ and containing every vertex. Distinct parametrizations of the same subspace count as one circle. The graph is uniquely Hamiltonian when it has exactly one such subspace.

The first question asks whether there exist an integer $r>2$ and a one-ended locally finite $r$-regular graph that is uniquely Hamiltonian. Here $r$-regular means that every vertex has degree $r$.

The source also asks whether such a uniquely Hamiltonian graph exists without the one-end restriction but with every vertex and every end having the same finite degree $r>2$. The degree of an end is the supremum of the sizes of families of pairwise vertex-disjoint rays belonging to it. These are separate existence questions; in the second, the degree requirement applies to all ends as well as to all vertices.
\subsection{Short English statement}
Can an infinite regular graph of degree greater than two have a unique Hamilton circle with only one end, or with every end having the same degree as its vertices?
\subsection{Sources}
\begin{itemize}
\item[S1] ULAM.AI and the UnsolvedMath Contributors, supplied problems.json, record 3299 (OPG-484), OpenGarden. Catalogue identity and source transcription; CC BY 4.0.
\url{https://huggingface.co/datasets/ulamai/UnsolvedMath}.
\item[S2] Open Problem Garden, Infinite uniquely hamiltonian graphs. Original problem and discussion.
\url{http://www.openproblemgarden.org/op/infinite_uniquely_hamiltonian_graphs}.
\item[S3] B. Mohar, Hamiltonicity of infinite graphs, Problem of the Month.
\url{https://www.fmf.uni-lj.si/~mohar/Problems/P0703_HamiltonicityInfinite.html}.
\end{itemize}
\end{document}
